\chapter[Analiza mikrolokalna]{Analiza mikrolokalna: zbiór falowy, symbole i operatory pseudoróżniczkowe}
\label{chap:pde31}

\section{Motywacja i konwencja transformaty Fouriera}
Nośnik osobliwy dystrybucji mówi, \emph{gdzie} dystrybucja
nie jest lokalnie reprezentowana przez funkcję gładką. Nie mówi jednak, w \emph{jakich kierunkach} jej
oscylacje wysokoczęstotliwościowe nie zanikają. Tę brakującą informację
koduje \emph{zbiór falowy} (ang. \emph{wavefront set}).
W analizie mikrolokalnej każdemu punktowi $x$ przyporządkowujemy
kowektor $\xi\ne0$; para $(x,\xi)$ należy do zbioru falowego, jeżeli
lokalna transformata Fouriera nie maleje szybko w stożku wokół $\xi$.

W tym rozdziale $d\ge1$, a $U\subset\R^d$ jest otwarty. Dla
$h\in\mathcal S(\R^d)$ ustalamy normalizację
\begin{equation}\label{eq:pde31-fourier}
 \widehat h(\xi)=\int_{\R^d}e^{-ix\cdot\xi}h(x)
 \,\mathrm d\lambda_d(x),\qquad
 h(x)=(2\pi)^{-d}\int_{\R^d}e^{ix\cdot\xi}\widehat h(\xi)
 \,\mathrm d\lambda_d(\xi).
\end{equation}
Przedłużamy ją na dystrybucje temperowane przez dualność,
zgodnie z definicją w rozdziale~\ref{chap:pde-fourier-Rd}.
Konwencja jest dokładnie tą z równania~\eqref{eq:pde8-convention};
nie zmieniamy znaku fazy ani czynników $(2\pi)^{-d}$.
Przez $\langle\xi\rangle=(1+|\xi|^2)^{1/2}$ oznaczamy standardową
wagę Sobolewa. \emph{Stożek otwarty} $\Gamma\subset\R^d\setminus\{0\}$
jest otwarty w $\R^d\setminus\{0\}$ i niezmienniczy
względem mnożenia przez dodatnie liczby.
Dystrybucje mogą być zespolone; stosujemy parowanie
liniowe względem funkcji testowej, zgodnie z konwencją
rozdziału~\ref{chap:pde-fourier-Rd}.

\begin{lemat}[Transformata dystrybucji o zwartym nośniku]
\label{lem:pde31-compact-fourier}
Niech $T\in\mathcal E'(\R^d)$ i niech
$\theta\in C_c^\infty(\R^d)$ będzie równa $1$
w otoczeniu $\operatorname{supp}T$. Wtedy
$\widehat T$ jest reprezentowana przez funkcję
\[
 F_T(\xi)=\langle T(x),\theta(x)e^{-ix\cdot\xi}\rangle.
\]
Funkcja ta nie zależy od wyboru $\theta$, jest gładka
i istnieje $M\in\N_0$ takie, że dla każdego
wielowskaźnika $\alpha$ istnieje $C_\alpha>0$ z
\[
 |\partial_\xi^\alpha F_T(\xi)|\le C_\alpha
 \langle\xi\rangle^M.
\]
\end{lemat}
\begin{proof}
Różnica dwóch dopuszczalnych funkcji $\theta$
znika w otoczeniu nośnika $T$, więc nie zmienia
wartości parowania. Z oszacowania lokalnego rzędu
dystrybucji (Twierdzenie~\ref{thm:pde7-local-order})
istnieją $C,M$ takie, że dla testów o nośniku
w ustalonym zwartym otoczeniu $\operatorname{supp}\theta$
wartość $|T(\psi)|$ nie przekracza
$C\max_{|\beta|\le M}\|\partial^\beta\psi\|_\infty$.
Stąd również $T$ przedłuża się na $\mathcal S$
przez $T(\psi)=T(\theta\psi)$, ciągle względem
półnorm Schwartza.
Różniczkując funkcję $\theta(x)e^{-ix\cdot\xi}$
względem $\xi$ w topologii funkcji testowych,
otrzymujemy
\[
 \partial_\xi^\alpha F_T(\xi)
 =\langle T(x),(-ix)^\alpha\theta(x)e^{-ix\cdot\xi}\rangle.
\]
Różniczkowanie jest dopuszczalne: na ustalonym zwartym
nośniku różnicowe ilorazy i wszystkie ich pochodne
względem $x$ zbiegają jednostajnie, z rozwinięcia
Taylora funkcji wykładniczej.
Reguła Leibniza i ograniczoność $x$ na nośniku
$\theta$ dają żądane oszacowanie, z tym samym $M$.

Dla $\psi\in\mathcal S$ całka z
$\theta(x)e^{-ix\cdot\xi}\psi(\xi)$ względem $\xi$
zbiega w każdej półnormie funkcji testowych
o ustalonym nośniku: pochodne względem $x$
są majorowane przez stałą razy
$\langle\xi\rangle^{|\beta|}|\psi(\xi)|$.
Można więc przenieść $T$ pod całkę
(przez obcięcie całki i przybliżenie sumami Riemanna).
Otrzymujemy
\[
 \int_{\R^d}F_T(\xi)\psi(\xi)\,\mathrm d\lambda_d(\xi)
 =\langle T,\theta\widehat\psi\rangle
 =\langle T,\widehat\psi\rangle
 =\langle\widehat T,\psi\rangle.
\]
To dowodzi reprezentacji transformaty.
\end{proof}

\begin{definicja}[Regularność mikrolokalna]\label{def:pde31-wf}
Niech $u\in\mathcal D'(U)$, $x_0\in U$ i
$\xi_0\in\R^d\setminus\{0\}$. Mówimy, że
$(x_0,\xi_0)\notin\operatorname{WF}(u)$, jeżeli istnieją
$\chi\in C_c^\infty(U)$ z $\chi=1$ w otoczeniu $x_0$
(iloczyn $\chi u$ przedłużamy przez zero na $\R^d$)
i otwarty stożek $\Gamma$ zawierający $\xi_0$, takie że dla każdego
$N\in\N$ istnieje $C_N>0$ spełniająca
\begin{equation}\label{eq:pde31-wf-decay}
 |\widehat{\chi u}(\xi)|\le C_N\langle\xi\rangle^{-N}
 \qquad(\xi\in\Gamma).
\end{equation}
Dopełnienie zbioru tak określonych par nazywamy zbiorem falowym $u$.
Jest to domknięty podzbiór $U\times(\R^d\setminus\{0\})$,
niezmienniczy względem dodatniego skalowania drugiej współrzędnej.
\end{definicja}
Istotnie, jeśli para jest regularna mikrolokalnie,
to ta sama funkcja $\chi$ i ten sam stożek świadczą
o regularności wszystkich par z iloczynu otoczenia,
na którym $\chi=1$, i stożka $\Gamma$.
Dopełnienie zbioru falowego jest więc otwarte.
Niezmienniczość wynika bezpośrednio ze stożkowości
$\Gamma$.

\begin{lemat}[Niezależność od dalszego obcięcia]\label{lem:pde31-cutoff}
Załóżmy, że $u\in\mathcal D'(U)$, $\chi\in C_c^\infty(U)$,
$\chi=1$ w otoczeniu $x_0$, i że
$\widehat{\chi u}$ szybko maleje w stożku $\Gamma$.
Jeżeli $\eta\in C_c^\infty(U)$ ma nośnik w tym otoczeniu,
a $\Gamma'$ jest stożkiem, którego domknięcie przecięte ze sferą
jednostkową leży wewnątrz $\Gamma$, to
$\widehat{\eta u}$ szybko maleje na $\Gamma'$.
\end{lemat}
\begin{proof}
Połóżmy $T=\chi u$, przedłużone przez zero na $\R^d$.
Ponieważ $\eta u=\eta T$, mamy
\[
 \widehat{\eta u}(\xi)=(2\pi)^{-d}\int_{\R^d}
 \widehat\eta(\xi-\zeta)\widehat T(\zeta)
 \,\mathrm d\lambda_d(\zeta).
\]
Istotnie, stosując odwrotną transformatę Fouriera
do $\eta$ i mnożąc przez $e^{-ix\cdot\xi}$,
otrzymujemy całkę funkcji wykładniczych z wagą
$\widehat\eta$. Po pomnożeniu przez funkcję
obcinającą równą $1$ blisko nośnika $T$,
całka zbiega w półnormach funkcji testowych,
jak w Lemacie~\ref{lem:pde31-compact-fourier}.
Możemy więc przenieść $T$ pod całkę i podstawić
$\zeta=\xi-\omega$, gdzie $\omega$ jest zmienną
całkowania we wzorze odwrotnym dla $\eta$.
Całka jest bezwzględnie zbieżna: z tego samego
lematu mamy $|\widehat T(\zeta)|\le C\langle\zeta\rangle^M$,
a $\widehat\eta$ należy do $\mathcal S$.

Z założenia o $\Gamma'$ istnieje $0<\delta<1/2$
takie, że dla $\xi\in\Gamma'$ i
$|\xi-\zeta|\le\delta|\xi|$ zachodzi $\zeta\in\Gamma$.
Wynika to z dodatniej odległości zwartego zbioru
$\overline{\Gamma'}\cap S^{d-1}$ od dopełnienia
$\Gamma\cap S^{d-1}$ na sferze: jednostkowe
kierunki $\xi$ i $\zeta$ różnią się o co najwyżej
$2|\xi-\zeta|/|\xi|\le2\delta$.
Jeśli dopełnienie na sferze jest puste,
wystarcza dowolne $0<\delta<1/2$.
W tym obszarze $|\zeta|\ge(1-\delta)|\xi|$,
więc dla dowolnego $N\in\N$
\[
 \int_{|\xi-\zeta|\le\delta|\xi|}
 |\widehat\eta(\xi-\zeta)\widehat T(\zeta)|
 \,\mathrm d\lambda_d(\zeta)
 \le C_N'\langle\xi\rangle^{-N}\|\widehat\eta\|_1 .
\]
W pozostałym obszarze kładziemy $w=\xi-\zeta$.
Mamy $|w|>\delta|\xi|$, a zatem
$\langle\xi\rangle\le\delta^{-1}\langle w\rangle$
oraz $\langle\xi-w\rangle\le(1+\delta^{-1})\langle w\rangle$.
Z wielomianowego wzrostu $\widehat T$ i szybkiego
zaniku $\widehat\eta$, wybierając wykładnik
$L=M+N+d+1$, dostajemy
\[
 \begin{aligned}
 &\int_{|w|>\delta|\xi|}
 |\widehat\eta(w)\widehat T(\xi-w)|
 \,\mathrm d\lambda_d(w)\\
 &\qquad\le C_L\int_{|w|>\delta|\xi|}
 \langle w\rangle^{-N-d-1}\,\mathrm d\lambda_d(w)\\
 &\qquad\le C_{L,\delta}\langle\xi\rangle^{-N}
 \int_{\R^d}\langle w\rangle^{-d-1}\,\mathrm d\lambda_d(w).
 \end{aligned}
\]
Ostatnia całka jest skończona, co widać
współrzędnymi biegunowymi dla $|w|\ge1$
i z ograniczoności funkcji na kuli jednostkowej.
Oba oszacowania dają szybki zanik w $\Gamma'$.
\end{proof}

\begin{stwierdzenie}[Rzut zbioru falowego]\label{prop:pde31-singular-support}
Dla $u\in\mathcal D'(U)$ zachodzi
$\operatorname{sing\,supp}(u)=\pi_x(\operatorname{WF}(u))$,
gdzie $\pi_x(x,\xi)=x$, a $\operatorname{sing\,supp}(u)$
oznacza zbiór punktów, w których $u$ nie jest funkcją klasy $C^\infty$
lokalnie.
\end{stwierdzenie}
\begin{proof}
Jeżeli $u$ jest gładka w otoczeniu $x_0$, to dla
$\chi\in C_c^\infty$ równej $1$ blisko $x_0$
i o nośniku w tym otoczeniu mamy
$\chi u\in C_c^\infty$; jej transformata maleje szybko we wszystkich
kierunkach, zatem $x_0$ nie leży w rzucie zbioru falowego.

Przeciwnie, załóżmy, że wszystkie pary $(x_0,\xi)$, $|\xi|=1$,
nie należą do zbioru falowego. Dla każdego kierunku $\xi$ wybieramy stożek
regularności $\Gamma_\xi$, funkcję obcinającą $\chi_\xi$
oraz mniejszy otwarty stożek $\Gamma_\xi'$
zawierający $\xi$, którego domknięcie na sferze
mieści się w $\Gamma_\xi$.
Ze zwartości sfery wybieramy skończone pokrycie
przez $\Gamma_{\xi_j}'$. W przecięciu otoczeń,
na których wszystkie $\chi_{\xi_j}=1$,
wybieramy $\eta\in C_c^\infty(U)$, równą $1$
blisko $x_0$. Z Lematu~\ref{lem:pde31-cutoff}
$\widehat{\eta u}$ szybko maleje w każdym
$\Gamma_{\xi_j}'$, a biorąc maksimum skończenie
wielu stałych, otrzymujemy szybki zanik na całej
$\R^d$ (również w zerze, przez ciągłość transformaty).
Definiujemy
\[
 v(x)=(2\pi)^{-d}\int_{\R^d}
 e^{ix\cdot\xi}\widehat{\eta u}(\xi)\,\mathrm d\lambda_d(\xi).
\]
Dla każdego wielowskaźnika $\alpha$ funkcja
$|\xi|^{|\alpha|}|\widehat{\eta u}(\xi)|$ jest całkowalna,
więc twierdzenie o zbieżności zmajoryzowanej
uzasadnia różniczkowanie pod całką i ciągłość
wszystkich pochodnych $v$.
Twierdzenie~\ref{thm:pde8-dist-rules} o odwracalności
transformaty w $\mathcal S'$ i twierdzenie Fubiniego
dla parowania z funkcjami Schwartza dają $v=\eta u$
w sensie dystrybucji.
Ponieważ $\eta=1$ w otoczeniu $x_0$, dystrybucja
$u$ jest tam reprezentowana przez funkcję gładką $v$.
\end{proof}

\section{Przykłady kierunkowych osobliwości}
\begin{przyklad}[Delta Diraca i skok]\label{ex:pde31-delta}
Dla $\delta_0\in\mathcal D'(\R^d)$ z Lematu
~\ref{lem:pde31-compact-fourier} mamy
$\widehat{\chi\delta_0}(\xi)=\chi(0)$.
Jeśli $\chi=1$ blisko zera, transformata jest
identycznie równa $1$ i nie maleje szybko w żadnym
niepustym stożku. Poza zerem dystrybucja jest
reprezentowana lokalnie przez funkcję zerową.
Zatem $\operatorname{WF}(\delta_0)=\{(0,\xi):\xi\ne0\}$.

Dla funkcji Heaviside'a $H(x)=\chi_{(0,\infty)}(x)$
w $\R$ zbiór falowy jest taki sam jak dla delty
w wymiarze jeden. Istotnie, dla dowolnej
$\chi\in C_c^\infty(\R)$ równej $1$ blisko zera
oraz $\xi\ne0$, całkowanie przez części daje
\[
 \widehat{\chi H}(\xi)
 =\int_0^\infty e^{-ix\xi}\chi(x)\,\mathrm d\lambda_1(x)
 =\frac{1}{i\xi}+
 \frac{1}{i\xi}\int_0^\infty e^{-ix\xi}\chi'(x)
 \,\mathrm d\lambda_1(x).
\]
Funkcja $\chi'\chi_{(0,\infty)}$ jest gładka
i ma zwarty nośnik, ponieważ $\chi'=0$
w otoczeniu zera. Jej transformata należy do
$\mathcal S(\R)$, więc ostatni składnik jest
$O(|\xi|^{-N})$ dla każdego $N$ przy $|\xi|\to\infty$.
W szczególności $|\xi\widehat{\chi H}(\xi)|\to1$
w obu kierunkach, co wyklucza szybki zanik
dla każdej dopuszczalnej lokalizacji.
Poza zerem $H$ jest lokalnie stała.
\end{przyklad}

\begin{przyklad}[Osobliwość hiperpłaszczyzny]\label{ex:pde31-conormal}
Niech $d\ge2$, $x=(x_1,x')$ i
$u=\delta(x_1)\otimes1(x')$, tj.
$\langle u,\varphi\rangle=\int_{\R^{d-1}}\varphi(0,x')
\,\mathrm d\lambda_{d-1}(x')$ dla $\varphi\in C_c^\infty(\R^d)$.
Wtedy
\[
 \operatorname{WF}(u)=
 \{((0,x'),(\xi_1,0)):\ x'\in\R^{d-1},\ \xi_1\ne0\}.
\]
Poza hiperpłaszczyzną $x_1=0$ dystrybucja znika
lokalnie. Ustalmy punkt $(0,x_0')$ na niej
i dowolną funkcję obcinającą $\chi$ równą $1$
blisko tego punktu. Z definicji dystrybucji
\[
 \widehat{\chi u}(\xi_1,\xi')
 =\int_{\R^{d-1}} e^{-ix'\cdot\xi'}\chi(0,x')
 \,\mathrm d\lambda_{d-1}(x')
 =\widehat h(\xi'),\qquad h(x')=\chi(0,x').
\]
Mamy $h\in C_c^\infty(\R^{d-1})$ i $h=1$
blisko $x_0'$. Jeśli kierunek $\xi_0=(\xi_{0,1},\xi_0')$
spełnia $\xi_0'\ne0$, to w odpowiednim otwartym
stożku wokół niego zachodzi $|\xi'|\ge c|\xi|$
z $c>0$. Szybki zanik $\widehat h$ daje więc
regularność mikrolokalną w tym stożku.
Jeśli natomiast $\xi_0'=0$, to z $h\not\equiv0$
i odwracalności transformaty Fouriera wynika,
że istnieje $\eta'\in\R^{d-1}$ z
$\widehat h(\eta')\ne0$. Dla $r\to+\infty$
wektory $(r\xi_{0,1},\eta')$ ostatecznie należą
do każdego otwartego stożka zawierającego $\xi_0$,
a transformata ma na nich tę samą niezerową wartość.
To wyklucza szybki zanik dla dowolnej $\chi$.
W szczególności nie wystarcza badanie samej półprostej
$\xi'=0$, gdyż całka z obcięcia $h$ może być zerowa.
\end{przyklad}

\begin{center}
\begin{tikzpicture}[scale=0.83,>=Latex]
  \draw[->,gray!70] (-2.8,0)--(2.8,0) node[right] {$x_1$};
  \draw[->,gray!70] (0,-2.5)--(0,2.6) node[above] {$x_2$};
  \draw[blue!65!black,very thick] (0,-2.25)--(0,2.25);
  \foreach \y in {-1.5,0,1.5}{
    \draw[-{Latex[length=3mm]},red!65!black,thick] (0,\y)--(1.25,\y);
    \draw[-{Latex[length=3mm]},red!65!black,thick] (0,\y)--(-1.25,\y);
  }
  \node[blue!65!black,below right] at (0,-2.25) {nośnik $u$};
  \node[red!65!black,above right] at (1.15,1.5) {kowektory normalne};
\end{tikzpicture}

\small Rysunek: nośnik $\delta(x_1)$ i kierunki zbioru falowego
w wymiarze dwa (schemat).
\end{center}

\begin{uwaga}[Zbiór falowy a nośnik]
Dystrybucja $\delta(x_1)$ ma nośnik będący hiperpłaszczyzną,
lecz zbiór falowy składa się z \emph{kowektorów normalnych} do niej.
Każda funkcja gładka, nawet o niezwartym nośniku,
ma pusty zbiór falowy; przykładem jest funkcja $1$.
Nie należy utożsamiać nośnika, nośnika osobliwego
ani zbioru falowego.
\end{uwaga}


\begin{twierdzenie}[Dokładna propagacja zbioru falowego przy translacji]
\label{thm:pde31-transport-wf}
Niech $u_0\in\mathcal D'(\R^d)$, $c\in\R^d$,
a dla $t\in\R$ niech $u(t)=\tau_{ct}u_0$,
gdzie $\tau_a u_0$ oznacza translację dystrybucji,
która dla funkcji ma postać $(\tau_a u_0)(x)=u_0(x-a)$.
Wtedy $t\mapsto u(t)$ jest gładkim odwzorowaniem
$\R\to\mathcal D'(\R^d)$ i rozwiązuje
$\partial_tu+c\cdot\nabla_xu=0$ przy $u(0)=u_0$.
Dla każdego $t\in\R$ zachodzi dokładna tożsamość
\begin{equation}\label{eq:pde31-transport-wf}
 \operatorname{WF}(u(t))=
 \{(x+ct,\xi):(x,\xi)\in\operatorname{WF}(u_0)\}.
\end{equation}
\end{twierdzenie}
\begin{proof}
Definiujemy translację przez
$\langle\tau_a u_0,\varphi\rangle
=\langle u_0,\varphi(\cdot+a)\rangle$.
Dla ustalonej funkcji testowej $\varphi$
i $t$ z przedziału zwartego $I$ funkcje
$\varphi(\cdot+ct)$ mają nośniki w jednym
zbiorze zwartym $K$. Taylorowski wzór z resztą
całkową pokazuje, że ilorazy różnicowe
zbiegają w każdej półnormie funkcji testowych
o nośniku w $K$. Dla pochodnych kolejnych rzędów
stosujemy ten sam argument do
$(c\cdot\nabla)^k\varphi(\cdot+ct)$.
Z ciągłości $u_0$ otrzymujemy
\[
 \frac{\mathrm d^k}{\mathrm dt^k}
 \langle u(t),\varphi\rangle
 =\langle u_0,(c\cdot\nabla)^k\varphi(\cdot+ct)\rangle
 =\langle(-c\cdot\nabla_x)^k u(t),\varphi\rangle .
\]
Zbieżność ilorazów różnicowych zachodzi także
w silnej topologii dystrybucji: dla ograniczonego
zbioru testów ich nośniki leżą w jednym zbiorze
zwartym i wszystkie półnormy są jednostajnie
ograniczone. Przesunięcia dla $t\in I$
zachowują te własności. Oszacowanie lokalnego
rzędu $u_0$ ogranicza resztę Taylora
po sparowaniu przez stałą razy moduł przyrostu,
jednostajnie na tym zbiorze testów.
Stąd $t\mapsto u(t)$ jest gładkie i
$\partial_tu=-c\cdot\nabla_xu$, a $u(0)=u_0$.

Dla dowolnej funkcji odcinającej
$\chi\in C_c^\infty(\R^d)$ definiujemy
$\eta(y)=\chi(y+a)$. Wtedy
$\chi\tau_a u_0=\tau_a(\eta u_0)$: dla testu
$\varphi$ obie strony mają wartość
$\langle u_0,\chi(\cdot+a)\varphi(\cdot+a)\rangle$.
Ponieważ $\eta u_0$ ma zwarty nośnik, wzór
z Lematu~\ref{lem:pde31-compact-fourier} i definicja
translacji dają
\[
 \widehat{\chi\tau_a u_0}(\xi)
 =e^{-ia\cdot\xi}\widehat{\eta u_0}(\xi).
\]
Współczynnik wykładniczy ma moduł jeden, zatem
szybki zanik w dowolnym stożku jest równoważny
po obu stronach. Funkcja $\chi$ jest równa jeden
w otoczeniu $x_0+a$ dokładnie wtedy, gdy $\eta$
jest równa jeden blisko $x_0$.
Stąd $(x_0+a,\xi_0)\notin\operatorname{WF}(\tau_a u_0)$
wtedy i tylko wtedy, gdy
$(x_0,\xi_0)\notin\operatorname{WF}(u_0)$.
Podstawiamy $a=ct$.
\end{proof}

\begin{przyklad}[Transport punktowej osobliwości]
\label{ex:pde31-moving-delta}
W jednym wymiarze dla $u_0=\delta_0$ otrzymujemy
$u(t,x)=\delta(x-ct)$. Z
\eqref{eq:pde31-transport-wf} wynika
$\operatorname{WF}(u(t))=\{(ct,\xi):\xi\ne0\}$.
Na czasoprzestrzeni $\R^2$ definiujemy $u$ przez
$\langle u,\varphi\rangle=\int_\R\varphi(t,ct)
\,\mathrm d\lambda_1(t)$.
Jest to dystrybucja rzędu zero: dla testów o nośniku
w ustalonym zbiorze zwartym całkowanie odbywa się
po ograniczonym przedziale czasu, więc moduł
parowania ogranicza stała razy $\|\varphi\|_\infty$.
Poza prostą $x=ct$ dystrybucja lokalnie znika.
Obliczymy
\[
 \operatorname{WF}(u)=
 \{((t,ct),(-c\xi,\xi)):t\in\R,\ \xi\ne0\}.
\]
Aby to sprawdzić, dla $\rho\in C_c^\infty(\R^2)$
obliczamy
$\widehat{\rho u}(\tau,\xi)=
\int_\R e^{-it(\tau+c\xi)}\rho(t,ct)\,
\mathrm d\lambda_1(t)$.
Jest to jednowymiarowa transformata funkcji
$h(t)=\rho(t,ct)\in C_c^\infty(\R)$ w punkcie
$\tau+c\xi$. W stożku wokół kierunku
$(\tau_0,\xi_0)$ z $\tau_0+c\xi_0\ne0$
mamy $|\tau+c\xi|\ge c_0|(\tau,\xi)|$
dla pewnego $c_0>0$, więc transformata
maleje szybko.
Ustalmy teraz niezerowy kierunek normalny
$(-c\xi_0,\xi_0)$ oraz dowolne obcięcie $\rho$
równe $1$ w otoczeniu $(t_0,ct_0)$.
Wtedy $h=1$ blisko $t_0$, więc z odwracalności
transformaty istnieje $\sigma\in\R$ z
$\widehat h(\sigma)\ne0$.
Dla $r\to+\infty$ pary
$(\sigma-rc\xi_0,r\xi_0)$ należą ostatecznie
do każdego otwartego stożka wokół kierunku
normalnego, ale transformata na nich jest
równa $\widehat h(\sigma)$. Wyklucza to
regularność mikrolokalną w tym kierunku.
Nie dowodzi to ogólnego twierdzenia o propagacji
bicharakterystyk dla operatorów o zmiennych współczynnikach.
\end{przyklad}


\begin{stwierdzenie}[Wygładzanie cieplne usuwa zbiór falowy]
\label{prop:pde-v44-heat-wavefront}
Niech $T\in\mathcal E'(\R^d)$ będzie dystrybucją
o zwartym nośniku i $t>0$. Dla jądra ciepła
$G_t(x)=(4\pi t)^{-d/2}e^{-|x|^2/(4t)}$
funkcja $G_t*T$ należy do $\mathcal S(\R^d)$,
a więc $\operatorname{WF}(G_t*T)=\varnothing$.
W szczególności dla $T=\delta_0$ zbiór falowy
w chwili początkowej jest niepusty we wszystkich
kierunkach, ale znika dla każdego $t>0$.
\end{stwierdzenie}
\begin{proof}
Z Lematu~\ref{lem:pde31-compact-fourier}
$\widehat T$ jest funkcją gładką na $\R^d$,
a ona sama i wszystkie jej pochodne mają
co najwyżej wzrost wielomianowy.
Jądro $G_t$ należy do $\mathcal S$, a jego
transformata wynosi $e^{-t|\xi|^2}$
(por. Lemat~\ref{lem:pde8-gauss}).
Twierdzenie~\ref{thm:pde8-convol-temp} daje
$\widehat{G_t*T}(\xi)=e^{-t|\xi|^2}\widehat T(\xi)$.
Każda pochodna mnożnika gaussowskiego jest
iloczynem wielomianu i tego samego mnożnika.
Z reguły Leibniza dla dowolnych $\alpha,N$
funkcja $\langle\xi\rangle^N
\partial_\xi^\alpha(e^{-t|\xi|^2}\widehat T(\xi))$
jest ograniczona, gdyż wykładniczy zanik
przeważa dowolną potęgę.
Iloczyn należy więc do przestrzeni Schwartza,
a jego odwrotna transformata również.
Z odwracalności transformaty w $\mathcal S'$
jest ona równa $G_t*T$ w sensie dystrybucji,
a ponieważ obie funkcje są gładkie,
równość zachodzi punktowo.
Z definicji~\ref{def:pde31-wf} każda funkcja gładka
ma pusty zbiór falowy.
\end{proof}
\begin{uwaga}[Nieodwracalność informacji wysokoczęstotliwościowej]
Powyższy fakt opisuje kierunki osobliwości, ale nie
jest oszacowaniem stabilności problemu odwrotnego.
Na obszarze ograniczonym mnożnik cieplny
$e^{-t\lambda_k}$ na funkcjach własnych laplasjanu
powoduje silną niestabilność odtwarzania warunku
początkowego; por. stwierdzenie
\ref{prop:pde-v44-heat-inverse} w rozdziale~\ref{chap:pde32}.
\end{uwaga}


\section{Symbole i kwantyzacja Kohna--Nirenberga}
\begin{definicja}[Klasa symboli $S^m_{1,0}$]\label{def:pde31-symbol}
Dla $m\in\R$ klasą $S^m_{1,0}(\R^d\times\R^d)$
nazywamy zbiór funkcji $a\in C^\infty(\R^d\times\R^d)$,
dla których dla każdych wielowskaźników $\alpha,\beta$
istnieje $C_{\alpha\beta}$ takie, że
\begin{equation}\label{eq:pde31-symbol}
 |\partial_x^\alpha\partial_\xi^\beta a(x,\xi)|
 \le C_{\alpha\beta}\langle\xi\rangle^{m-|\beta|}
 \quad(x,\xi\in\R^d).
\end{equation}
Dla $u\in\mathcal S(\R^d)$ operator pseudoróżniczkowy
w kwantyzacji Kohna--Nirenberga określamy wzorem
\begin{equation}\label{eq:pde31-op}
 \operatorname{Op}(a)u(x)=(2\pi)^{-d}
 \int_{\R^d}e^{ix\cdot\xi}a(x,\xi)\widehat u(\xi)
 \,\mathrm d\lambda_d(\xi).
\end{equation}
Liczba $m$ jest górnym ograniczeniem rzędu symbolu:
dla $m_1\le m_2$ mamy $S^{m_1}_{1,0}\subset S^{m_2}_{1,0}$.
Nie zakładamy jednorodności symbolu w $\xi$.
Całka~\eqref{eq:pde31-op} jest bezwzględnie zbieżna,
ponieważ $|a(x,\xi)|\le C_{00}\langle\xi\rangle^m$
i $\widehat u\in\mathcal S$.
\end{definicja}

\begin{lemat}[Ciągłość na przestrzeni Schwartza]
\label{lem:pde31-schwartz-op}
Jeżeli $a\in S^m_{1,0}$, to $\operatorname{Op}(a)$
jest ciągłym operatorem liniowym
$\mathcal S(\R^d)\to\mathcal S(\R^d)$.
\end{lemat}
\begin{proof}
Dla wielowskaźnika $\alpha$ różniczkujemy pod całką:
\[
 \partial_x^\alpha\operatorname{Op}(a)u(x)
 =(2\pi)^{-d}\int_{\R^d}e^{ix\cdot\xi}
 b_\alpha(x,\xi)\widehat u(\xi)\,\mathrm d\lambda_d(\xi),
 \qquad
 b_\alpha(x,\xi)=
 \sum_{\gamma\le\alpha}\binom{\alpha}{\gamma}
 (i\xi)^{\alpha-\gamma}\partial_x^\gamma a(x,\xi).
\]
Z reguły Leibniza i oszacowań symbolu wynika
$|\partial_\xi^\nu b_\alpha(x,\xi)|
\le C_{\alpha\nu}\langle\xi\rangle^{m+|\alpha|-|\nu|}$.
Każdy iloczyn $b_\alpha\widehat u$ jest więc
jednostajnie względem $x$ szybko malejący wraz
ze wszystkimi pochodnymi po $\xi$.
W szczególności funkcje pod całkami otrzymanymi
przez kolejne różniczkowania względem $x$
mają całkowalne majoranty niezależne od $x$.
Uzasadnia to różniczkowanie przez twierdzenie
o zbieżności zmajoryzowanej.

Dla wielowskaźnika $\beta$ wykorzystujemy
$x^\beta e^{ix\cdot\xi}=i^{-|\beta|}
\partial_\xi^\beta e^{ix\cdot\xi}$.
Całkowanie przez części daje
\[
 x^\beta\partial_x^\alpha\operatorname{Op}(a)u(x)
 =(2\pi)^{-d}i^{|\beta|}
 \int_{\R^d}e^{ix\cdot\xi}
 \partial_\xi^\beta(b_\alpha(x,\xi)\widehat u(\xi))
 \,\mathrm d\lambda_d(\xi).
\]
Nie ma wyrazów brzegowych dzięki szybkiemu
zanikowi iloczynu i jego pochodnych.
Rozwijając pochodną przez regułę Leibniza,
dla liczby całkowitej $L> d+\max(0,m+|\alpha|)$
otrzymujemy
\[
 \sup_x|x^\beta\partial_x^\alpha\operatorname{Op}(a)u(x)|
 \le C_{\alpha\beta L}
 \sum_{\nu\le\beta}\sup_\xi
 \langle\xi\rangle^L|\partial_\xi^\nu\widehat u(\xi)|.
\]
Istotnie, po wyłączeniu wskazanych półnorm
pozostaje całkowalna majoranta
$\langle\xi\rangle^{m+|\alpha|-L}$.
Prawa strona jest ciągłą półnormą funkcji $u$
z ciągłości transformaty Fouriera na przestrzeni
Schwartza (rozdział~\ref{chap:pde-fourier-Rd}).
Półnormy po lewej wyznaczają jej topologię,
co dowodzi tezy.
\end{proof}

\begin{przyklad}[Mnożenie, pochodne, potencjały Bessela]
Jeżeli $b\in C_b^\infty(\R^d)$ (wszystkie pochodne są ograniczone),
to $a(x,\xi)=b(x)$ należy do $S^0_{1,0}$ i
$\operatorname{Op}(a)u=bu$.
Dla $a(\xi)=i\xi_j\in S^1_{1,0}$ dostajemy
$\partial_{x_j}$ przez różniczkowanie wzoru
odwrotnego Fouriera. Ogólniej, dla
$P=\sum_{|\alpha|\le k}b_\alpha(x)\partial^\alpha$
z $b_\alpha\in C_b^\infty$ symbolem jest
$p(x,\xi)=\sum_{|\alpha|\le k}b_\alpha(x)(i\xi)^\alpha
\in S^k_{1,0}$; kolejność mnożenia i różniczkowania
jest dokładnie taka jak we wzorze na $P$.

Symbol $\langle\xi\rangle^m$ definiuje operator
$(I-\Delta)^{m/2}$ (potencjał Bessela dla $m<0$).
Przynależność do $S^m_{1,0}$ wynika z następującego
oszacowania: każda pochodna rzędu $|\beta|$
jest sumą wyrazów postaci
$p(\xi)(1+|\xi|^2)^{m/2-j}$, gdzie $p$
jest wielomianem stopnia co najwyżej $2j-|\beta|$.
Dowodzimy tego indukcyjnie: różniczkowanie $p$
zmniejsza stopień o jeden, a różniczkowanie potęgi
zwiększa $j$ o jeden i mnoży $p$ przez współrzędną.
Każdy taki wyraz jest ograniczony przez stałą
razy $\langle\xi\rangle^{m-|\beta|}$.
Wzór na operator wynika z definicji mnożnika
Fouriera; por. Twierdzenie~\ref{thm:pde8-Hs}.
\end{przyklad}

\begin{twierdzenie}[Oszacowanie Sobolewa dla symboli o zwartym nośniku w $x$]
\label{thm:pde31-compact-symbol}
Niech $m,s\in\R$ i $a\in S^m_{1,0}$ ma nośnik w zmiennej $x$
zawarty w ustalonym zwartym zbiorze $K\subset\R^d$.
Wtedy $\operatorname{Op}(a)$ przedłuża się jednoznacznie do
ograniczonego operatora $H^s(\R^d)\to H^{s-m}(\R^d)$.
Jego norma zależy od $m,s,d,K$ i skończonej liczby
stałych z~\eqref{eq:pde31-symbol}.
\end{twierdzenie}
\begin{proof}
Najpierw niech $u\in\mathcal S(\R^d)$.
Niech $\widehat a_x(\eta,\xi)$ oznacza transformatę
Fouriera $a$ po zmiennej $x$. Dla każdej liczby
całkowitej $k\ge0$ mamy
\[
 \langle\eta\rangle^{2k}\widehat a_x(\eta,\xi)
 =\int_{\R^d}e^{-ix\cdot\eta}(I-\Delta_x)^ka(x,\xi)
 \,\mathrm d\lambda_d(x).
\]
Całkowanie przez części nie daje wyrazów brzegowych,
bo $a(\cdot,\xi)$ ma nośnik w $K$.
Moduł prawej strony jest nie większy niż
$C_k\lambda_d(K)\langle\xi\rangle^m$,
gdzie $C_k$ zależy tylko od stałych symbolu
dla pochodnych po $x$ rzędu co najwyżej $2k$.
W szczególności dla dowolnego $N\ge0$,
wybierając $2k\ge N$, otrzymujemy
$|\widehat a_x(\eta,\xi)|\le
C_N\langle\eta\rangle^{-N}\langle\xi\rangle^m$.

Z twierdzenia Fubiniego
\[
 \widehat{\operatorname{Op}(a)u}(\zeta)
 =(2\pi)^{-d}\int_{\R^d}
 \widehat a_x(\zeta-\xi,\xi)\widehat u(\xi)
 \,\mathrm d\lambda_d(\xi).
\]
Istotnie, całka z modułu funkcji pod podwójną
całką nie przekracza
$C_{00}\lambda_d(K)\int_{\R^d}
\langle\xi\rangle^m|\widehat u(\xi)|\,\mathrm d\lambda_d(\xi)$,
co jest skończone.

Dla $s\in\R$ kładziemy
$v(\xi)=\langle\xi\rangle^s\widehat u(\xi)$.
Jądro operatora przekształcającego $v$ w
$\langle\zeta\rangle^{s-m}\widehat{\operatorname{Op}(a)u}(\zeta)$
wynosi
\[
 J(\zeta,\xi)=(2\pi)^{-d}\langle\zeta\rangle^{s-m}
 \widehat a_x(\zeta-\xi,\xi)\langle\xi\rangle^{-s}.
\]
Nierówność $\langle x+y\rangle\le
\sqrt2\langle x\rangle\langle y\rangle$
wynika z $|x+y|^2\le2|x|^2+2|y|^2$.
Stosując ją do $\zeta=\xi+(\zeta-\xi)$
dla $s-m\ge0$ albo do
$\xi=\zeta+(\xi-\zeta)$ dla $s-m<0$,
dostajemy nierówność Peetre'a
\[
 \langle\zeta\rangle^{s-m}\langle\xi\rangle^{m-s}
 \le2^{|s-m|/2}\langle\zeta-\xi\rangle^{|s-m|}.
\]
Stąd $|J(\zeta,\xi)|\le
C_N\langle\zeta-\xi\rangle^{-N+|s-m|}$.
Wybieramy $N>d+|s-m|$ i oznaczamy
\[
 A_1=\sup_\zeta\int_{\R^d}|J(\zeta,\xi)|
 \,\mathrm d\lambda_d(\xi),\qquad
 A_2=\sup_\xi\int_{\R^d}|J(\zeta,\xi)|
 \,\mathrm d\lambda_d(\zeta).
\]
Obie liczby są skończone przez podstawienie
$\eta=\zeta-\xi$ i całkowalność
$\langle\eta\rangle^{-N+|s-m|}$.
Nierówność Schura wynika tutaj bezpośrednio
z nierówności Cauchy'ego--Schwarza:
\[
 \left|\int_{\R^d}J(\zeta,\xi)v(\xi)
 \,\mathrm d\lambda_d(\xi)\right|^2
 \le A_1\int_{\R^d}|J(\zeta,\xi)|\,|v(\xi)|^2
 \,\mathrm d\lambda_d(\xi).
\]
Całkując po $\zeta$ i stosując twierdzenie Tonellego,
otrzymujemy oszacowanie kwadratu normy $L^2$
przez $A_1A_2\|v\|_2^2$.
W normach Sobolewa po obu stronach występuje
ten sam czynnik $(2\pi)^{-d/2}$, zgodny z
normalizacją Plancherela. W konsekwencji
\[
 \|\operatorname{Op}(a)u\|_{H^{s-m}}
 \le\sqrt{A_1A_2}\|u\|_{H^s}.
\]
Z gęstości $\mathcal S$ w $H^s$ i zupełności
$H^{s-m}$ (Twierdzenie~\ref{thm:pde8-Hs})
wynika istnienie i jednoznaczność ciągłego
przedłużenia. Podane oszacowania pokazują
zależność jego normy tylko od wymienionych danych.
\end{proof}

\begin{uwaga}[Zasięg twierdzenia]
Zwartego nośnika symbolu w zmiennej $x$ nie wymaga pełne
twierdzenie o ograniczoności operatorów pseudoróżniczkowych,
lecz nasze założenie umożliwia samodzielny dowód przez jądro Fouriera.
Dla symboli ogólnych potrzebna jest dodatkowa lokalizacja
lub pełna teoria ograniczoności operatorów pseudoróżniczkowych;
nie utożsamiamy tego skrótu z jej dowodem.
\end{uwaga}

\begin{wniosek}[Lokalne oszacowanie operatora]
\label{cor:pde31-local-sobolev}
Dla $a\in S^m_{1,0}$, $\rho\in C_c^\infty(\R^d)$
i $s\in\R$ operator $\rho\operatorname{Op}(a)$
ma jednoznaczne ograniczone przedłużenie
$H^s(\R^d)\to H^{s-m}(\R^d)$.
\end{wniosek}
\begin{proof}
Na przestrzeni Schwartza
$\rho\operatorname{Op}(a)=\operatorname{Op}(\rho a)$
bez żadnego wyrazu resztowego, bezpośrednio
ze wzoru~\eqref{eq:pde31-op}.
Z reguły Leibniza $\rho a\in S^m_{1,0}$:
pochodne $\rho$ są ograniczone, a pochodne
po $\xi$ działają jedynie na $a$.
Nośnik w zmiennej $x$ mieści się w
$\operatorname{supp}\rho$. Stosujemy
Twierdzenie~\ref{thm:pde31-compact-symbol}.
\end{proof}

\section{Jądro poza przekątną i lokalność osobliwości}
\begin{stwierdzenie}[Gładkość jądra poza przekątną]
\label{prop:pde31-kernel}
Niech $a\in S^m_{1,0}$. Jądro dystrybucyjne
$K_a(x,y)$ operatora~\eqref{eq:pde31-op} jest funkcją gładką
na zbiorze $\{(x,y):x\ne y\}$.
\end{stwierdzenie}
\begin{proof}
Formalnie
$K_a(x,y)=(2\pi)^{-d}\int_{\R^d}
 e^{i(x-y)\cdot\xi}a(x,\xi)\,\mathrm d\lambda_d(\xi)$.
Najpierw wstawiamy obcięcie częstotliwości
$\rho(\xi/R)\in C_c^\infty$ równe jeden dla $|\xi|\le R$.
Dla $z=x-y\ne0$ zachodzi
$e^{iz\cdot\xi}=(i|z|^2)^{-1}z\cdot\nabla_\xi
 e^{iz\cdot\xi}$. Po $N$ całkowaniach przez części
pochodne symbolu maleją jak $\langle\xi\rangle^{m-N}$.
Na każdym zwartym podzbiorze $\{|x-y|>0\}$ wybieramy
$N>m+d+k$ dla pochodnych do rzędu $k$, co zapewnia
jednostajną zbieżność całek i ich pochodnych przy $R\to\infty$.
Ponieważ $k$ jest dowolne, jądro jest gładkie poza przekątną.
\end{proof}


\begin{lemat}[Przedłużenie działania na dystrybucje o zwartym nośniku]
\label{lem:pde31-distribution-op}
Niech $a\in S^m_{1,0}(\R^d\times\R^d)$, $m\in\R$.
Operator $\operatorname{Op}(a)$ określony wzorem
\eqref{eq:pde31-op} dla funkcji Schwartza ma
kanoniczne przedłużenie
$\mathcal E'(\R^d)\to\mathcal D'(\R^d)$.
Dla $u\in\mathcal E'$ i $\varphi\in C_c^\infty$ jest ono
zdefiniowane przez
\[
\langle\operatorname{Op}(a)u,\varphi\rangle
=\langle u,\operatorname{Op}(a)^t\varphi\rangle,
\]
gdzie $t$ oznacza transpozycję względem biliniowego
parowania dystrybucji z funkcją testową.
\end{lemat}
\begin{proof}
Dla $\varphi\in C_c^\infty(\R^d)$ formalna transpozycja ma
postać
\[
 (\operatorname{Op}(a)^t\varphi)(y)
 =(2\pi)^{-d}\int_{\R^d}e^{-iy\cdot\xi}
 \left(\int_{\R^d}e^{ix\cdot\xi}a(x,\xi)\varphi(x)
 \,\mathrm d\lambda_d(x)\right)\mathrm d\lambda_d(\xi).
\]
Całka wewnętrzna jest szybko malejąca jako funkcja $\xi$.
W istocie, na zwartym nośniku $\varphi$ pochodne w $x$
symbolu $a$ są oszacowane przez $C\langle\xi\rangle^m$.
Wielokrotne całkowanie przez części względem $x$ daje dla
każdego $N$ oszacowanie $C_N\langle\xi\rangle^{m-N}$.
Ten sam argument działa po różniczkowaniu w $\xi$,
a zatem całka wewnętrzna należy do $\mathcal S(\R^d)$.
Jej odwrotna transformata jest funkcją gładką w $y$,
łącznie ze wszystkimi pochodnymi. Ponieważ $u$ ma zwarty
nośnik, działa w sposób ciągły na funkcjach gładkich
w otoczeniu swojego nośnika. Powyższy wzór definiuje
zatem dystrybucję $\operatorname{Op}(a)u$; ciągłość względem
$\varphi$ wynika z oszacowań na skończoną liczbę
półnorm funkcji testowych. Dla $u\in C_c^\infty$ dostajemy
pierwotną całkę przez twierdzenie Fubiniego.
\end{proof}

\begin{wniosek}[Własność pseudolokalna]
\label{cor:pde31-pseudolocal}
Jeśli $a\in S^m_{1,0}$, a $u\in\mathcal E'(\R^d)$
jest dystrybucją o zwartym nośniku, przy czym
$\operatorname{Op}(a)u$ rozumiemy w sensie
lematu~\ref{lem:pde31-distribution-op}, to
$\operatorname{sing\,supp}(\operatorname{Op}(a)u)
\subset\operatorname{sing\,supp}(u)$.
\end{wniosek}
\begin{proof}
Niech $x_0$ leży poza zbiorem osobliwości $u$.
Rozkładamy $u=\rho u+(1-\rho)u$, gdzie $\rho=1$ blisko $x_0$
i jest wsparta w otoczeniu gładkości $u$.
Funkcja $\rho u$ jest gładka i ma zwarty nośnik,
więc jej transformata jest szybko malejąca, a
\eqref{eq:pde31-op} można różniczkować dowolnie wiele razy.
Drugie wyrażenie daje funkcję gładką przy $x_0$ dzięki
stwierdzeniu~\ref{prop:pde31-kernel}, gdyż nośnik
$(1-\rho)u$ nie przecina otoczenia $x_0$.
\end{proof}

\section{Parametryks i regularność eliptyczna}
\begin{twierdzenie}[Parametryks operatora o stałych współczynnikach]
\label{thm:pde31-parametrix}
Niech $P(D)=\sum_{|\alpha|\le m}c_\alpha\partial^\alpha$
będzie operatorem o stałych współczynnikach o rzędzie $m\ge1$.
Załóżmy, że jego jednorodny symbol główny
$P_m(i\xi)=\sum_{|\alpha|=m}c_\alpha(i\xi)^\alpha$
nie znika dla $\xi\ne0$. Wtedy istnieje
$q\in S^{-m}_{1,0}$ i funkcja $r\in C_c^\infty(\R^d)$
takie, że
\begin{equation}\label{eq:pde31-param}
 q(\xi)P(i\xi)=1-r(\xi).
\end{equation}
Dla każdej dystrybucji temperowanej $u$ zachodzi
$u=q(D)P(D)u+r(D)u$. Operator $r(D)$ jest wygładzający,
tj. przekształca $\mathcal S'$ w $C^\infty$.
Jeśli $P(D)u\in H^s(\R^d)$ oraz $u\in H^t(\R^d)$ dla pewnego $t\in\R$,
to $u\in H^{s+m}(\R^d)$.
\end{twierdzenie}
\begin{proof}
Z eliptyczności symbolu głównego i zwartości sfery wynika
$|P_m(i\xi)|\ge c|\xi|^m$. Ponieważ symbole niższych
rzędów są $O(|\xi|^{m-1})$, mamy
$|P(i\xi)|\ge c'\langle\xi\rangle^m$ przy $|\xi|\ge R$.
Wybieramy $\rho\in C^\infty$ równą zero dla $|\xi|\le R$
i jeden dla $|\xi|\ge2R$. Kładziemy
$q(\xi)=\rho(\xi)/P(i\xi)$ i $r=1-\rho$.
Wielokrotne różniczkowanie ilorazu pokazuje
$|\partial_\xi^\beta q(\xi)|\le C_\beta
\langle\xi\rangle^{-m-|\beta|}$, a więc $q\in S^{-m}_{1,0}$.
Otrzymujemy~\eqref{eq:pde31-param}.
Ponieważ $r$ ma zwarty nośnik w częstotliwościach,
$\partial^\alpha r(D)u$ jest funkcją gładką dla każdego
wielowskaźnika $\alpha$, co wynika z reprezentacji Fouriera
dystrybucji o zwartym nośniku w przestrzeni częstotliwości.
Wreszcie $q(D):H^s\to H^{s+m}$ jest ograniczony
z definicji norm Sobolewa, a $r(D):H^t\to H^{s+m}$
dla dowolnych $s,t$, gdyż mnożnik ma zwarty nośnik i jest gładki.
\end{proof}

\begin{uwaga}[Eliptyczność mikrolokalna]
Dla gładkich operatorów o zmiennych współczynnikach analogiczny
parametryks konstruuje się iteracyjnie w rachunku symboli.
Standardowa konsekwencja to
$\operatorname{WF}(Pu)\subset\operatorname{WF}(u)$ oraz
$\operatorname{WF}(u)\subset\operatorname{WF}(Pu)\cup\operatorname{Char}(P)$,
gdzie $\operatorname{Char}(P)$ jest zbiorem zer symbolu głównego
poza przekrojem zerowym. Pełny dowód ogólnego rachunku składania
symboli wykracza poza przedstawiony tu samodzielny dowód
parametryksu dla stałych współczynników; nie stosujemy tego
wyniku do operatorów z jedynie mierzalnymi współczynnikami.
\end{uwaga}

\section{Zbiór falowy Sobolewa i powierzchnie charakterystyczne}
\begin{definicja}[Zbiór falowy Sobolewa rzędu $s$]
\label{def:pde31-wfs}
Dla $s\in\R$, $u\in\mathcal D'(U)$, para $(x_0,\xi_0)$
nie należy do $\operatorname{WF}^s(u)$, jeśli istnieje
$\chi\in C_c^\infty(U)$ równa jeden blisko $x_0$ i stożek
$\Gamma\ni\xi_0$ taki, że
$\int_\Gamma\langle\xi\rangle^{2s}|\widehat{\chi u}(\xi)|^2
\,\mathrm d\lambda_d(\xi)<\infty$.
\end{definicja}

\begin{przyklad}[Próg Sobolewa delty]
Dla $\delta_0$ lokalna transformata jest stała, więc
$\delta_0\in H^s_{\mathrm{loc}}(\R^d)$ dokładnie wtedy, gdy
$s<-d/2$. W istocie w współrzędnych biegunowych
zbieżność całki jest równoważna warunkowi
$\int_1^\infty r^{2s+d-1}\,\mathrm d\lambda_1(r)<\infty$.
Dla $s\ge-d/2$ cały stożek kierunków w punkcie zero
należy do $\operatorname{WF}^s(\delta_0)$.
\end{przyklad}

\begin{przyklad}[Charakterystyki równania falowego]
Symbol główny $\partial_t^2-c^2\Delta_x$ wynosi
$-\tau^2+c^2|\xi|^2$. Jego zbiór charakterystyczny jest stożkiem
$\tau^2=c^2|\xi|^2$, $(\tau,\xi)\ne0$.
W odróżnieniu od laplasjanu, który nie ma rzeczywistych
niezerowych kowektorów charakterystycznych, operator falowy
może przenosić osobliwości wzdłuż charakterystyk.
Twierdzenie o propagacji zbioru falowego wzdłuż bicharakterystyk
jest głębsze niż uzyskana tutaj własność pseudolokalna:
nie przedstawiamy jego pełnego dowodu bez rachunku operatorów
całkowych Fouriera.
\end{przyklad}

\section{Przybliżenie optyki geometrycznej}
\begin{stwierdzenie}[Równanie eikonalne jako warunek wysokich częstotliwości]
\label{prop:pde31-eikonal}
Niech $c>0$, a $a,\varphi$ będą rzeczywistymi funkcjami klasy
$C^2$ na otwartym podzbiorze $\R^{1+d}$.
Dla $u_\eps=a\exp(i\varphi/\eps)$ zachodzi
\[
 (\partial_t^2-c^2\Delta_x)u_\eps
 =e^{i\varphi/\eps}
 \Bigl[\eps^{-2}a(c^2|\nabla_x\varphi|^2-\varphi_t^2)
 +i\eps^{-1}\mathcal T(a,\varphi)
 +(\partial_t^2-c^2\Delta_x)a\Bigr],
\]
gdzie
$\mathcal T(a,\varphi)=2\varphi_ta_t+a\varphi_{tt}
 -2c^2\nabla_x\varphi\cdot\nabla_xa-c^2a\Delta_x\varphi$.
Zanik wyrazu rzędu $\eps^{-2}$ w punktach, gdzie $a\ne0$,
wymaga równania eikonalnego $\varphi_t^2=c^2|\nabla_x\varphi|^2$.
\end{stwierdzenie}
\begin{proof}
Różniczkujemy dwukrotnie iloczyn $ae^{i\varphi/\eps}$:
każda pochodna wykładnika wnosi $i\partial\varphi/\eps$.
Po zebraniu współczynników przy $\eps^{-2}$ otrzymujemy
$a(c^2|\nabla_x\varphi|^2-\varphi_t^2)$,
a przy $i\eps^{-1}$ dokładnie $\mathcal T(a,\varphi)$.
Pozostałe wyrazy odpowiadają operatorowi falowemu
zastosowanemu do amplitudy $a$.
\end{proof}

\section{Jak stosować analizę mikrolokalną w PDE?}
Praktyczny schemat polega na identyfikacji symbolu głównego,
wyznaczeniu zbioru charakterystycznego i zbadaniu,
czy rozważane pary $(x,\xi)$ są eliptyczne.
Tam, gdzie symbol nie znika, parametryks usuwa możliwość
występowania nowych osobliwości. Na zbiorze charakterystycznym
istotna staje się geometria przepływu hamiltonowskiego symbolu.
Przed użyciem tego schematu trzeba określić \emph{gładkość
współczynników}; metody z tego rozdziału nie zastępują teorii
De Giorgiego--Nasha--Mosera dla macierzy mierzalnych.



